\begin{proof}[Proof of \cref{obj:prop:unit-overlap-boundary}]
Write
\[
\alpha:=\exp(-1/t_0),
\qquad
L:=\exp(\zeta).
\]
Since \(t_0>0\) and \(\zeta>0\), one has \(0<\alpha<1\) and \(1<L\). % lean: CausalSmith.Stat.PomdpLatentOverlapMinimax.unit_overlap_boundary

1. Define
\[
c_{\mathrm{imm}}
:=
2\left\{
L\left(1+\frac{4}{L-1}\right)+\frac{4}{1-\alpha}
\right\}.
\]
The preceding inequalities imply \(c_{\mathrm{imm}}>0\). For any \(T\ge1\), \cref{obj:lem:radius-sensitive-phiw} applied at depth \(k=0\) and overlap constant \(C=1\) gives, uniformly over \(M\in\mathcal M_T(t_0,\zeta,1)\),
\[
\mathbb E_M\!\left[
\bigl(\widehat\theta_0-\theta(K,e)\bigr)^2
\right]
\le
\frac{2}{T}
\left\{
L\left(1+\frac{4}{L-1}\right)+\frac{4}{1-\alpha}
\right\}
=
\frac{c_{\mathrm{imm}}}{T}.
\]
Taking the supremum over the model class yields
\[
\sup_{M\in\mathcal M_T(t_0,\zeta,1)}
\mathbb E_M\!\left[
\bigl(\widehat\theta_0-\theta(K,e)\bigr)^2
\right]
\le
\frac{c_{\mathrm{imm}}}{T}.
\]
% lean: CausalSmith.Stat.PomdpLatentOverlapMinimax.radius_sensitive_phiw

2. Let \(M\in\mathcal M_T(t_0,\zeta,1)\). Write \(d_b\) and \(d_e\) for the behavior and target stationary laws on the finite joint state space \(\mathcal S\), and define
\[
\left\|d_b-d_e\right\|_{\mathrm{TV}}
:=
\frac12\sum_{s\in\mathcal S}|d_b(s)-d_e(s)| .
\]
The stationary-law total-variation bound at overlap radius \(q_C\) gives
\[
\left\|d_b-d_e\right\|_{\mathrm{TV}}
\le q_1,
\]
and at unit overlap \cref{obj:def:overlap-distance-radius} gives
\[
q_1=\frac{1-1}{1}=0.
\]
Thus
\[
\left\|d_b-d_e\right\|_{\mathrm{TV}}
\le0.
\]
Since the displayed total-variation quantity is a half-sum of absolute values, it is nonnegative, so
\[
\left\|d_b-d_e\right\|_{\mathrm{TV}}=0.
\]
Unfolding the definition of total variation yields
\[
\sum_{s\in\mathcal S}|d_b(s)-d_e(s)|=0.
\]
For each joint state \(s\), the nonnegative summand is bounded by the whole finite sum,
\[
0\le |d_b(s)-d_e(s)|
\le
\sum_{u\in\mathcal S}|d_b(u)-d_e(u)|
=0,
\]
and hence
\[
|d_b(s)-d_e(s)|=0.
\]
Therefore \(d_b(s)=d_e(s)\) for every \(s\in\mathcal S\), and consequently \(d_e=d_b\). % lean: CausalSmith.Stat.PomdpLatentOverlapMinimax.stationary_tv_le_overlapRadius

3. Set
\[
c_{\mathrm{lo}}:=\frac{\exp(-1/4)}{64},
\qquad
c_{\mathrm{hi}}:=\max\{c_{\mathrm{imm}},c_{\mathrm{lo}}\},
\qquad
T_\star:=1.
\]
Then \(c_{\mathrm{lo}}>0\), \(c_{\mathrm{lo}}\le c_{\mathrm{hi}}\), and therefore \(c_{\mathrm{hi}}>0\); also \(T_\star\ge1\). For every \(T\ge T_\star\), \cref{obj:lem:uniform-parametric-floor} at \(C=1\) gives
\[
\frac{\exp(-1/4)}{64T}
\le
R_T(t_0,\zeta,1),
\]
which is
\[
\frac{c_{\mathrm{lo}}}{T}
\le
R_T(t_0,\zeta,1).
\]
For the upper bound, the minimax risk is bounded above by the worst-case risk of any admissible estimator with nonnegative squared risk; applying this to \(\widehat\theta_0\) and using Step 1 gives
\[
R_T(t_0,\zeta,1)
\le
\sup_{M\in\mathcal M_T(t_0,\zeta,1)}
\mathbb E_M\!\left[
\bigl(\widehat\theta_0-\theta(K,e)\bigr)^2
\right]
\le
\frac{c_{\mathrm{imm}}}{T}
\le
\frac{c_{\mathrm{hi}}}{T}.
\]
% lean: CausalSmith.Stat.PomdpLatentOverlapMinimax.uniform_parametric_floor

4. Finally,
\[
t_0\zeta>0,
\qquad\text{hence}\qquad
2+t_0\zeta>2.
\]
Therefore
\[
\frac{2}{2+t_0\zeta}<1,
\]
which is the asserted bound on the hidden-state exponent. % lean: CausalSmith.Stat.PomdpLatentOverlapMinimax.unit_overlap_boundary
\end{proof}
